% Nêu cụ thể mục tiêu của KLTN, đặc biệt phải nêu được mục tiêu cải tiến sẽ là gì so với thực trạng nêu trong phần tổng quan đề tài

\subsection*{a) Mục tiêu tổng quát:}

Mục tiêu của đề tài là xây dựng một kiến trúc MLOps cho các mô hình học sâu đa phương thức, nhằm hỗ trợ quản lý và tự động hóa vòng đời của mô hình một cách hiệu quả. Hệ thống được định hướng theo nền tảng hiện đại trên nền tảng điện toán đám mây, bảo đảm khả năng mở rộng, tính ổn định và vận hành bền vững.


\subsection*{b) Mục tiêu cụ thể:}
\begin{itemize}
    \item \textbf{Thu thập và chuẩn hóa bộ dữ liệu đa nguồn:} 
    Xây dựng quy trình thu thập, tiền xử lý và chuẩn hóa dữ liệu từ nhiều nguồn dữ liệu khác nhau (ảnh, văn bản, v.v.). Đảm bảo dữ liệu được lưu trữ có cấu trúc, có gắn nhãn, kiểm soát chất lượng và dễ dàng truy xuất phục vụ cho quá trình huấn luyện và đánh giá mô hình.

    \item \textbf{Xây dựng và tự động hóa quy trình huấn luyện mô hình:} 
    Thiết kế pipeline MLOps cho quá trình huấn luyện, đánh giá và triển khai mô hình học máy đa phương thức. Tích hợp các công cụ tự động hóa giúp giảm thao tác thủ công, đảm bảo khả năng tái lập thí nghiệm và rút ngắn thời gian phát triển mô hình.

    \item \textbf{Thiết lập cơ chế giám sát và tái huấn luyện:} 
    Xây dựng hệ thống theo dõi hiệu suất mô hình sau khi triển khai, phát hiện suy giảm chất lượng (model drift, data drift) và kích hoạt quy trình tái huấn luyện khi cần thiết. Đảm bảo hệ thống hoạt động ổn định và duy trì hiệu suất theo thời gian.

    \item \textbf{Đánh giá và kiểm chứng kết quả mô hình:} 
    Xây dựng quy trình đánh giá mô hình dựa trên các chỉ số phù hợp với bài toán, thực hiện kiểm thử trên các tập dữ liệu độc lập và so sánh với các phương pháp hiện có. Đảm bảo kết quả mô hình có tính tin cậy, ổn định và đáp ứng yêu cầu ứng dụng thực tế.

\end{itemize}
